% =============================================================================
\section{Connections to NLP Context Evaluation}\label{app:nlp}
% =============================================================================

This appendix formally maps between genomic ECL and the NLP effective context length literature.

\subsection{NIAH as a special case of influence energy}

The Needle-in-a-Haystack (NIAH) test inserts a specific fact (the ``needle'') at position $d$ in a distractor context (the ``haystack'') and tests whether the model retrieves it. In our framework, this corresponds to a specific perturbation kernel and influence energy:

Let $X_{\mathrm{hay}}$ be a haystack sequence and $X_{\mathrm{hay}}^{+\mathrm{needle}@d}$ be the sequence with the needle inserted at position $d$. The NIAH sensitivity is:
\[
\mathrm{NIAH}(d) = \E\!\left[d_{\mathrm{task}}\!\left(g(f(X_{\mathrm{hay}})),\, g(f(X_{\mathrm{hay}}^{+\mathrm{needle}@d}))\right)\right],
\]
which is exactly a task-aware influence energy $\cI^{\mathrm{task}}(d; r_{\mathrm{query}})$ under a specific ``insertion'' perturbation rather than a ``replacement'' perturbation. The NIAH-based ECL is $\max\{d : \mathrm{NIAH}(d) > \tau\}$, which is a special case of task-aware ECL$_\beta$ with a hard threshold $\tau$ replacing the fractional threshold $\beta$.

\subsection{RULER operationalization and its relation to perturbation ECL}

RULER~\citep{Hsieh2024RULER} defines ECL as:
\[
\ECL_{\mathrm{RULER}} = \max\{L : \mathrm{Score}(L) \ge \mathrm{Score}(L_{\mathrm{ref}})\},
\]
where $\mathrm{Score}(L)$ is average performance across 13 tasks at context length $L$. This is a \emph{task-specific, threshold-based} ECL that depends on the task suite.

Our perturbation-variance ECL is complementary:
\begin{itemize}[leftmargin=2em]
\item It operates at the embedding level (task-agnostic) or task-specific level.
\item It is a continuous, fractional measure rather than binary.
\item It provides statistical uncertainty quantification.
\item It does not require carefully designed ``needles'' or ``tasks''---any perturbation kernel suffices.
\end{itemize}

\subsection{LongPPL and key-token identification}

LongPPL~\citep{Fang2024LongPPL} identifies key tokens by comparing next-token predictions with and without long context:
\[
\mathrm{KeyScore}(i) = \log P(x_i \mid x_{<i}) - \log P(x_i \mid x_{i-w:i}),
\]
where $w$ is a short context window. Tokens with large positive KeyScore genuinely benefit from long-range context.

In our framework, KeyScore is approximately related to out-of-window ablation:
\[
\mathrm{KeyScore}(i) \approx -\Delta_{\mathrm{out}}(w; i)
\]
under a log-loss discrepancy. Key tokens are positions where $\Delta_{\mathrm{out}}$ is large, i.e., positions whose predictions degrade substantially when distal context is removed.

\subsection{Lost-in-the-middle as position-dependent ECL}

The U-shaped performance curve in ``Lost in the Middle''~\citep{Liu2024LostInMiddle} can be interpreted as position-dependent influence energy. If we define a ``query position'' $r$ at the end of the sequence (where the model generates its answer), the influence profile $\cI(i;r)$ should ideally be uniform across $i$ (all context positions equally useful). The U-shape reveals $\cI(i;r)$ is high for $i$ near the beginning and end but low for $i$ in the middle---a form of position-dependent ECL degradation.

In genomics, the analogue would be if models preferentially attend to positions near the edges of their input window while underweighting the center. This can be tested by plotting $\cI(d;r)$ for reference positions $r$ at various locations within the input window.

\subsection{Intrinsic entropy and optimal context length}

\citet{Shi2025Entropy} proved that an optimal context length $L^\star$ exists, beyond which additional context increases approximation loss. In our framework, this corresponds to the regime where $\cI(d;r) < 0$ (negative influence), meaning that additional context \emph{hurts} predictions. This can occur when out-of-distribution perturbations or distant noise degrades the model. Our framework can detect this by checking whether $\Delta_{\mathrm{out}}(\ell;r)$ is non-monotone in $\ell$: if $\Delta_{\mathrm{out}}$ \emph{decreases} as $\ell$ grows (removing more context improves predictions), the model suffers from context poisoning.
